\begin{result}{E.2.22}{An exact reduction to one tower of covers}{Proved}
Put \(X=\mathbf P^1_{\Q}\setminus\{0,1,\infty\}\), and retain
the actual Betti lattices and comparison maps of E.2.21.
For \(\sigma\in\Gamma_{F_{\mathrm{cm}}}\), the following
conditions are equivalent:
\begin{enumerate}
\item \(\rho_A(\sigma)\Lambda_A=\Lambda_{{}^{\sigma}A}\)
for every abelian variety \(A/\Qb\).
\item For every connected finite \(\acute{e}\)tale cover
\(Y\to X_{\Qb}\), with smooth projective compactification
\(C_Y\), the actual Galois comparison on the Jacobian satisfies
\[
 \rho_{J(C_Y)}(\sigma)\Lambda_{J(C_Y)}
       =\Lambda_{J({}^{\sigma}C_Y)}.
 \tag{E.2.22a}
\]
\end{enumerate}
There is a precise topological formulation of condition 2.
Fix \(x=2\), let \(\Pi=\pi_1(X(\C),x)\), and identify
\(\widehat\Pi\) with \(\pi_1^{\mathrm{et}}(X_{\Qb},x)\).
The rational point gives the Galois action
\(\alpha_\sigma\in\Aut(\widehat\Pi)\).
For every open subgroup \(U\subset\widehat\Pi\), set
\[
 \Pi_U=\Pi\cap U,\qquad
 \Lambda_U=H_1(C_U(\C),\Z),
 \tag{E.2.22b}
\]
where \(C_U\) compactifies the pointed cover associated to \(U\).
Then condition 2 says exactly that the induced map
\[
 \bar\alpha_{\sigma,U}:
 \Lambda_U\otimes\widehat\Z
       \longrightarrow
 \Lambda_{\alpha_\sigma(U)}\otimes\widehat\Z
 \tag{E.2.22c}
\]
carries \(\Lambda_U\) onto \(\Lambda_{\alpha_\sigma(U)}\)
for every \(U\). This property depends only on the outer
Galois action. Consequently these conditions describe exactly
the subgroup \(\Sigma_F\) in E.2.21.
\end{result}
\begin{proof}
We first supply the integral quotient statement needed for the
converse. Every nonzero abelian variety \(A/\Qb\) admits a
smooth projective curve \(C/\Qb\) and a homomorphism
\[
 q:J(C)\longrightarrow A,\qquad
 q_*H_1(J(C)(\C),\Z)=H_1(A(\C),\Z).
 \tag{E.2.22d}
\]
For dimension one take \(C=A\). In higher dimension take the
successive smooth ample sections in the proof of
\cite[Theorem 10.1, Lemma 10.3, pp.33--35]{GenMilneJVs2021}.
For every finite surjective \(\acute{e}\)tale isogeny
\(A'\to A\), their successive inverse images are smooth and
connected: pullback preserves ampleness, an ample effective
divisor on a projective variety of dimension at least two is
connected, and smoothness is preserved by an \(\acute{e}\)tale
base change. Thus the resulting curve has connected inverse
image under every such isogeny. Its inclusion induces the
surjective \(q\) of that theorem.

Here is why this particular construction gives the asserted
integral surjectivity. Write
\(A(\C)=V/\Lambda_A\), and put
\(\Lambda'=q_*H_1(J(C)(\C),\Z)\).
Surjectivity of \(q\) makes \(\Lambda'\) a subgroup of finite
index in \(\Lambda_A\). Suppose that index exceeds one.
Choose \(n\geq1\) with \(n\Lambda_A\subset\Lambda'\).
The complex torus \(V/\Lambda'\), with its map to \(A\),
comes from an actual isogeny over \(\Qb\): form
\[
 H=(n^{-1}\Lambda')/\Lambda_A\subset A[n],
 \qquad A/H\longrightarrow A,
 \tag{E.2.22e}
\]
where the last map is induced by \([n]\).
Multiplication by \(n\) identifies its analytic source with
\(V/\Lambda'\). All points of \(A[n]\) and hence this finite
quotient are defined over \(\Qb\).
The image of \(\pi_1(C(\C))\) in
\(\pi_1(A(\C))=\Lambda_A\) is precisely \(\Lambda'\),
because the Abel--Jacobi map is an isomorphism on integral
first homology. The pullback of (E.2.22e) to \(C\) therefore
has trivial monodromy on \(\Lambda_A/\Lambda'\); it is a
disjoint union of more than one copy of \(C(\C)\).
This contradicts the connected inverse-image property.
It proves (E.2.22d). The same argument applies to
\({}^{\sigma}q\).

Belyi's construction, in the primary version
\cite[Theorem 1 and old proof, pp.1--2]{GenBelyi1997}, gives a
finite morphism \(C\to\mathbf P^1_{\Qb}\) whose branch
values lie in \(\{0,1,\infty\}\). Removing the inverse image
of these three points makes \(C\) the compactification of a
cover occurring in condition 2.
The relevant actual homology/Tate maps commute:
\[
\begin{array}{ccc}
 \Lambda_{J(C)}\otimes\widehat\Z
   &\xrightarrow{\rho_{J(C)}(\sigma)}&
 \Lambda_{J({}^{\sigma}C)}\otimes\widehat\Z\\
 \big\downarrow q_* && \big\downarrow({}^{\sigma}q)_*\\
 \Lambda_A\otimes\widehat\Z
   &\xrightarrow{\rho_A(\sigma)}&
 \Lambda_{{}^{\sigma}A}\otimes\widehat\Z.
\end{array}
 \tag{E.2.22f}
\]
In these completed modules, (E.2.22d) and condition 2 imply
the exact equality
\[
\begin{split}
 \rho_A(\sigma)\Lambda_A
  &=\rho_A(\sigma)q_*\Lambda_{J(C)}\\
  &=({}^{\sigma}q)_*\rho_{J(C)}(\sigma)\Lambda_{J(C)}\\
  &=({}^{\sigma}q)_*\Lambda_{J({}^{\sigma}C)}
    =\Lambda_{{}^{\sigma}A}.
\end{split}
 \tag{E.2.22g}
\]
The zero abelian variety contributes no constraint.
This proves 2 implies 1; 1 implies 2 by taking the indicated
Jacobians.

For completeness, the map in (E.2.22c) uses the compactified
homology quotient, not merely the homology of the affine
cover. The finite-index subgroup \(\Pi_U\) has profinite
completion \(U\). In its abelianization, quotient by the
classes of positively oriented loops around the punctures;
this gives \(\Lambda_U\). Equivalently, quotient
\(U^{\mathrm{ab}}\) by the closed subgroup generated by its
peripheral inertia images; the result is
\(\Lambda_U\otimes\widehat\Z\).
The Galois action carries those inertia groups to the inertia
groups of the conjugate cover, and hence induces
\(\bar\alpha_{\sigma,U}\).

The Abel--Jacobi map induces the abelianized fundamental-group
isomorphism of \cite[Proposition 9.1 and Corollary 9.6,
pp.30--32]{GenMilneJVs2021}. Together with finite-coefficient
Betti--\(\acute{e}\)tale comparison
\cite[Expos\'e XVI, Th\'eor\`eme 4.1]{GenSGA4}, it gives
the commutative diagram
\[
\begin{array}{ccc}
 \Lambda_U\otimes\widehat\Z
   &\xrightarrow{\bar\alpha_{\sigma,U}}&
 \Lambda_{\alpha_\sigma(U)}\otimes\widehat\Z\\
 \big\downarrow\simeq && \big\downarrow\simeq\\
 T J(C_U)&\xrightarrow{\rho_{J(C_U)}(\sigma)}&
 T J(C_{\alpha_\sigma(U)}).
\end{array}
 \tag{E.2.22h}
\]
Here \(TJ=\prod_\ell T_\ell J\). On the Jacobian the
comparison sends a lattice vector \(v\bmod n\) to
\(v/n\bmod\Lambda_J\); on the curve it is induced by
Abel--Jacobi. Changing its base point is a translation on
the Jacobian and acts trivially on first homology.
Thus (E.2.22h) identifies the actual integral lattices as
well as their completions. Algebraic cover maps, their
Jacobian maps and deck actions commute with this diagram;
no independent choices of transport are introduced.

Finally, inner automorphisms do not change this lattice
condition. Given \(h\in\widehat\Pi\) and open \(U\),
density supplies \(p\in\Pi\) with \(h\in pU\).
Writing \(h=pu\), \(u\in U\), gives
\(hUh^{-1}=pUp^{-1}\), and conjugation by \(h\) on
\(U^{\mathrm{ab}}\) equals conjugation by \(p\), because
conjugation by \(u\) is inner. The latter is induced by an
actual topological change of base path, so it carries the
compactified integral homology lattice onto the corresponding
lattice. This proves the asserted outer independence.
\end{proof}

\paragraph{The remaining map.}
E.2.22 is a proved equivalence of the existing all-abelian
lattice condition with the stated condition on one cover
tower. It does not establish that its outer automorphisms lie
in the image of a profinite mapping-class group.
The opened primary theorem
\cite[Theorem 1.1, p.2, proof p.3]{GenMatsumotoTamagawa1997}
identifies the intersection of the Galois image and the
orientation-preserving profinite mapping-class image as the
identity for an affine hyperbolic curve with rational
punctures. The map from the integral condition (E.2.22c) to
that mapping-class image remains unconstructed here.
Completed cup or Reidemeister pairing preservation alone
does not supply it: the actual \(\acute{e}\)tale Galois maps
already preserve those pairings up to the cyclotomic
character, which is one on \(\Gamma_{F_{\mathrm{cm}}}\).
The quartic place-pair and the full pure-motive problem
therefore remain open.
